\documentclass[11pt]{article}
\usepackage[margin=1in]{geometry}
\usepackage[T1]{fontenc}
\usepackage{graphicx}
\usepackage{hyperref}
\usepackage{enumitem}
\setlength{\parskip}{0.5em}
\setlength{\parindent}{0pt}

\title{diffusion models for image generation}
\author{arXiv GraphRAG --- phase-6 deep research}
\date{\today}

\begin{document}
\maketitle

\begin{abstract}
The review is excellent on theory and speed-up methods but lacks coverage of the dominant transformer-based architectures and spatial control mechanisms central to the 2023+ frontier.
\end{abstract}

\section{The Evolution of Diffusion Models: From Foundational Generative Modeling to the Visual Frontier}

The central challenge of generative modeling lies in the ability to accurately capture and sample from high-dimensional, complex data distributions to produce realistic, novel imagery. For years, the field was dominated by paradigms such as Generative Adversarial Networks (GANs) and Variational Autoencoders (VAEs), which, despite their successes, often faced limitations regarding training stability, mode collapse, and sample diversity. The emergence of diffusion models has fundamentally shifted this landscape, offering a mathematically grounded framework that provides superior sample quality and more stable convergence. As these models have matured, they have transitioned from theoretical curiosities to the backbone of modern visual synthesis, driving a revolution in how machines interpret, manipulate, and create visual content.

This review traces the trajectory of this technological evolution, beginning with the foundations of generative modeling and the inception of early diffusion processes. We examine the subsequent mathematical and structural refinements, specifically the integration of score-based models and the move toward latent space representations to enhance computational efficiency. The discussion then moves to the semantic revolution brought about by text-to-image synthesis and the mechanisms of conditional guidance that enable precise user control. To address the practical bottlenecks of iterative refinement, we explore advancements in efficient sampling and model distillation, as well as the expansion of diffusion techniques into the domains of image restoration and low-level vision. Finally, we look toward the current frontier, investigating the cutting-edge developments in geometric, multi-scale, and training-free diffusion methods that define the next generation of generative AI.

\subsection{Foundations of Generative Modeling and Early Diffusion}

The evolution of generative modeling has been defined by a tension between the efficient, adversarial training of GANs and the more stable, likelihood-based approaches of variational and diffusion models. Early foundational work in latent variable modeling established the use of stochastic variational inference and reparameterization of the variational lower bound to enable efficient learning in directed probabilistic models with continuous latent variables \texttt{arXiv:1312.6114}. This approach was later contrasted by the introduction of Generative Adversarial Networks (GANs), which shifted the paradigm toward a minimax two-player game where a generator \$G\$ and discriminator \$D\$ are trained simultaneously via backpropagation to capture data distributions \texttt{arXiv:1406.2661}.

While GANs provided a framework for high-fidelity synthesis, they faced significant training challenges. Efforts to stabilize these models included the Two Time-Scale Update Rule (TTUR), which utilized individual learning rates for the discriminator and generator to encourage convergence to a local Nash equilibrium \texttt{arXiv:1706.08500}. Further architectural refinements, such as style-based generators, sought to improve distribution quality and disentanglement by separating high-level attributes from stochastic variations via scale-specific control \texttt{arXiv:1812.04948}.

Parallel to the adversarial lineage, a different trajectory emerged through non-equilibrium thermodynamics. This approach utilized an iterative forward diffusion process to destroy data structure, followed by a learned reverse diffusion process to achieve tractability in modeling complex datasets \texttt{arXiv:1503.03585}. This line of inquiry was significantly advanced by score-based generative models, which estimated gradients of the data distribution via score matching on data perturbed with varying Gaussian noise levels, using annealed Langevin dynamics for sampling \texttt{arXiv:1907.05600}.

The convergence of these ideas culminated in the development of Denoising Diffusion Probabilistic Models (DDPMs) and score-based modeling via Stochastic Differential Equations (SDEs). DDPMs achieved high-quality image synthesis by training on a weighted variational bound that connects diffusion processes to denoising score matching \texttt{arXiv:2006.11239}. This was further unified by the framework of SDEs, which provides a continuous-time generalization of both DDPMs and score-based modeling by transforming data distributions into priors via forward SDEs and reversing the process through a reverse-time SDE \texttt{arXiv:2011.13456}. This SDE-based approach offers significant novelty by introducing a probability flow ODE for exact likelihood and deterministic sampling, achieving state-of-the-art performance with an FID of 2.20 and an Inception Score of 9.89 on CIFAR-10 \texttt{arXiv:2011.13456}. However, despite these mathematical strengths, the framework remains computationally demanding, as sampling is significantly slower than GAN-based methods and requires complex hyperparameter tuning for SDE selection and sampler configuration \texttt{arXiv:2011.13456}.

\subsection{Advancing Diffusion: Score-based Models and Latent Spaces}

The evolution of diffusion models has been characterized by a dual focus on accelerating sampling processes and shifting generative modeling from pixel space to more efficient latent representations. Early advancements centered on optimizing the sampling trajectory. Denoising Diffusion Implicit Models (DDIMs) introduced non-Markovian diffusion processes that maintain the same training objective as DDPMs but facilitate significantly faster reverse sampling, achieving 10x to 50x speedups in wall-clock time [\texttt{arXiv:2010.02502}]. This acceleration was further pursued through knowledge distillation, where multi-step denoising processes are distilled into a single-step "denoising student" [\texttt{arXiv:2101.02388}]. While DDIMs [\texttt{arXiv:2010.02502}] and distillation methods [\texttt{arXiv:2101.02388}] prioritize speed, other works have focused on refining the underlying probabilistic framework. For instance, improving DDPMs by learning the variances of the reverse diffusion process allows for competitive log-likelihoods and enhanced sample quality [\texttt{arXiv:2102.09672}].

A parallel track of research has addressed the stability and scalability of score-based generative models. By analyzing learning in high-dimensional spaces and employing an exponential moving average (EMA) of model weights, researchers improved the ability to scale training to higher resolutions, such as 256x256 [\texttt{arXiv:2006.09011}]. This emphasis on quality and distribution coverage culminated in the demonstration that diffusion models could outperform GANs on ImageNet, utilizing classifier guidance to trade off diversity for fidelity [\texttt{arXiv:2105.05233}]. This capability is further refined by elucidating the design space of these models, separating choices in sampling, training, and preconditioning to achieve superior FID scores, such as 1.36 on ImageNet-64 [\texttt{arXiv:2206.00364}].

The most significant shift in computational efficiency involves moving the diffusion process into latent spaces. Latent Score-based Generative Models (LSGM) utilize a variational autoencoder framework and a novel score-matching objective to improve sampling speed and distribution coverage [\texttt{arXiv:2106.05931}]. Building upon this concept, Latent Diffusion Models (LDMs) apply diffusion within the latent space of pretrained autoencoders and introduce cross-attention layers to enable flexible conditioning on inputs like text or bounding boxes [\texttt{arXiv:2112.10752}]. While LSGM [\texttt{arXiv:2106.05931}] focuses on variance reduction and score function parameterization, LDMs [\texttt{arXiv:2112.10752}] prioritize reducing the overall computational requirements of training and inference while maintaining high quality. Together, these advancements represent a transition from computationally intensive pixel-level modeling to highly efficient, conditioned latent-space generation.

\subsection{Text-to-Image Synthesis and Conditional Guidance}

The evolution of text-to-image synthesis has transitioned from basic conditional generation to sophisticated hierarchical and language-intensive frameworks. Early approaches, such as GLIDE, focused on photorealistic generation and editing by comparing CLIP guidance against classifier-free guidance strategies [\texttt{arXiv:2112.10741}]. While GLIDE established a foundation for text-conditional diffusion, subsequent research sought to refine the relationship between text embeddings and visual synthesis. For instance, a hierarchical two-stage approach was introduced to improve image diversity and caption similarity by utilizing a prior to generate CLIP image embeddings, which then guide a diffusion decoder [\texttt{arXiv:2204.06125}].

A significant shift in the field involves the scaling of language understanding to improve image-text alignment. Rather than focusing solely on the scale of the image diffusion model, research has demonstrated that prioritizing large, pretrained text-only language models (such as T5) as encoders can significantly enhance sample fidelity [\texttt{arXiv:2205.11487}]. This approach achieved a FID score of 7.27 on the COCO dataset and demonstrated superior performance on the DrawBench benchmark [\texttt{arXiv:2205.11487}]. This emphasis on deep language understanding contrasts with earlier models that relied more heavily on CLIP-based embeddings [\texttt{arXiv:2112.10741}, \texttt{arXiv:2204.06125}].

Central to the control of these models is the mechanism of guidance. The introduction of classifier-free diffusion guidance revolutionized the field by jointly training conditional and unconditional diffusion models, allowing for a tunable trade-off between sample fidelity and diversity [\texttt{arXiv:2207.12598}]. While this method is highly effective, it necessitates joint training of both conditional and unconditional components [\texttt{arXiv:2207.12598}]. To address the computational overhead associated with such training-heavy methods, more recent work has proposed training-free guidance for pretrained unconditional models [\texttt{arXiv:2312.13253}]. This training-free alternative offers a significant practical advantage, reducing computation time by approximately threefold while maintaining comparable image quality [\texttt{arXiv:2312.13253}].

Finally, the scope of diffusion models has expanded from pure synthesis to high-fidelity image editing. While earlier models like GLIDE provided editing capabilities through fine-tuning for inpainting [\texttt{arXiv:2112.10741}], newer methodologies incorporate rectifier modules to modulate weights with residual features [\texttt{arXiv:2312.15707}]. This novel learning paradigm aims to minimize error propagation, facilitating more precise high-fidelity reconstruction and out-of-domain editing [\texttt{arXiv:2312.15707}]. Consequently, the literature shows a clear trajectory from establishing basic text-conditional guidance toward optimizing the efficiency of guidance and the depth of linguistic comprehension.

\subsection{Efficient Sampling and Model Distillation}

The primary bottleneck in diffusion models remains the high computational cost of iterative sampling. Research to address this has diverged into two main trajectories: training-free acceleration via advanced ODE solvers and model distillation techniques that reduce sampling steps through additional training.

Training-free methods focus on improving the efficiency of the underlying differential equations. DPM-Solver [\texttt{arXiv:2206.00927}] achieves significant speedups (4-16x) by formulating diffusion ODE solutions through analytical computation of the linear part and an exponentially weighted integral, guaranteeing convergence order. This approach allows for high-quality sampling in as few as 10 steps (FID 4.70 on CIFAR-10) without requiring retraining. Similarly, [\texttt{arXiv:2408.08306}] proposes accelerating the process by utilizing an exponentially time-decaying Signal to Noise Ratio and learning diffusion coefficients via an autoencoder, reducing steps from 1000 to 200–500.

In contrast, distillation-based methods aim for extreme acceleration, often targeting one-step generation, but they necessitate additional training phases. Progressive distillation [\texttt{arXiv:2202.00512}] iteratively halves the number of sampling steps, achieving an FID of 3.0 on CIFAR-10 in just 4 steps. This can be extended to guided models; [\texttt{arXiv:2210.03142}] employs a two-stage distillation process—matching combined conditional/unconditional outputs followed by progressive distillation—to achieve up to 256x speedup.

The pursuit of one-step generation has led to more specialized architectures. Consistency models [\texttt{arXiv:2303.01469}] directly map noise to data, enabling one-step generation with an FID of 3.55 on CIFAR-10, and can be trained either via distillation or as standalone models. Recent advancements have further refined this: [\texttt{arXiv:2506.09376}] utilizes lightweight GAN fine-tuning on pre-trained models (freezing 85\% of parameters) to convert them into one-step generators, while Denoising Score Distillation (DSD) [\texttt{arXiv:2503.07578}] focuses on distilling models from noisy or corrupted samples into high-quality one-step generators.

A notable tension exists between the complexity of training and the resulting inference speed. While DPM-Solver [\texttt{arXiv:2206.00927}] provides a training-free solution, it still requires multiple steps. Distillation methods like those in [\texttt{arXiv:2202.00512}] and [\texttt{arXiv:2210.03142}] offer much higher speedups but require pre-trained models and significant optimization. To bridge these needs, Score-based Idempotent Generative Networks (SIGNs) [\texttt{arXiv:2509.21470}] offer a middle ground; they distill idempotent models from diffusion scores to map distributions back to a target manifold, enabling a flexible quality-efficiency trade-off without the need for adversarial training. This allows for high-fidelity generation across various source domains, including corrupted distributions, while maintaining lower sampling costs.

\subsection{Diffusion for Image Restoration and Low-Level Vision}

The application of diffusion models to low-level vision tasks has evolved from specialized training strategies to training-free exploitation of pre-trained architectures. A primary challenge in this domain is the alignment of the diffusion process with specific degradation models. For instance, \texttt{arXiv:2302.02398} addresses this by redefining the diffusion process to account for Gaussian, Gamma, and Poisson noise models. By estimating the posterior distribution of clean images conditioned on these specific noise types, the model aims to achieve generative denoising that preserves fine details and avoids the over-smoothing common in traditional methods. This approach, however, necessitates a unified training strategy to handle diverse noise contexts.

In contrast to models that require task-specific training, \texttt{arXiv:2402.16907} introduces Diffusion Posterior Proximal Sampling (DPPS), a training-free framework that leverages pre-trained diffusion models. By utilizing measurement-consistent sample selection and measurement-signal initialization, DPPS seeks to enhance restoration performance and output stability when dealing with degraded images. This represents a significant shift toward efficiency, as it avoids the computational burden of retraining models for every new degradation type, though it remains dependent on the quality of the underlying pre-trained priors.

Further advancements have focused on the efficiency and structural fidelity of the diffusion process. \texttt{arXiv:2306.13720} proposes a method using analytical image attenuation to perform simultaneous image-to-zero and zero-to-noise mapping. This approach derives specific training objectives and sampling formulas that allow for high-quality unconditioned generation and various downstream tasks—such as super-resolution, inpainting, and edge detection—while significantly accelerating the process through fewer denoising steps. While this improves efficiency, the reliance on analytical attenuation functions defines the scope of its applicability.

Recent research has also targeted the preservation of fine structural details through noise scheduling. \texttt{arXiv:2410.01540} implements a hybrid noise scheme within diffusion and flow-matching frameworks, utilizing an edge-aware scheduler. This scheduler transitions from edge-preserving noise to isotropic noise, specifically designed to improve robustness and perceptual quality in structure-guided tasks like stroke-to-image synthesis. The effectiveness of this method is measured through FID, KID, and CLIP-score, highlighting a move toward more perceptually grounded evaluation.

Collectively, these works reflect a broader trend in the field, as summarized in the survey by \texttt{arXiv:2406.11138}, which categorizes the progress of diffusion models in low-level vision through the lenses of performance and efficiency. While early methods like \texttt{arXiv:2302.02398} focused on redefining the diffusion process for specific noise distributions, contemporary research has branched into training-free sampling (\texttt{arXiv:2402.16907}), accelerated analytical mapping (\texttt{arXiv:2306.13720}), and sophisticated edge-aware noise scheduling (\texttt{arXiv:2410.01540}).

\subsection{The Frontier: Geometric, Multi-Scale, and Training-Free Diffusion}

Recent advancements in diffusion models have shifted from fundamental denoising towards optimizing the underlying data manifold, multi-scale feature representations, and inference-time efficiency. A significant trend involves decomposing image signals to improve both speed and quality. MSF achieves a 4x speed-up by factorizing latent features into a low-frequency base and a high-frequency residual signal (\texttt{arXiv:2501.13349}). This multi-scale approach is echoed in the spectral domain by Spectral Progressive Diffusion, which utilizes an optimal resolution schedule derived from a model's power spectrum to enable efficient image and video generation (\texttt{arXiv:2605.18736}). While MSF focuses on latent factorization, Hierarchical Koopman Diffusion offers a more structural approach by lifting nonlinear dynamics into a latent space governed by linear Koopman operators across scale-specific subspaces, facilitating fast one-step sampling (\texttt{arXiv:2510.12220}).

The geometric structure of the data manifold also presents a critical frontier. The MIND model explicitly models manifold geometry by integrating discrete patch tokenization into a continuous score function via soft top-k aggregation and dual-branch high-frequency embedding (\texttt{arXiv:2606.00094}). This geometric awareness yields high fidelity, with MIND-XL achieving an FID of 1.95 under guidance (\texttt{arXiv:2606.00094}). Similarly, efforts to improve convergence and quality include Gaussianization preprocessing to make training distributions resemble independent Gaussians (\texttt{arXiv:2512.21020}), and the use of truncated Karhunen-Loève expansion to represent Brownian motion, which accelerates training convergence (\texttt{arXiv:2503.17657}).

A distinct divergence exists between models requiring heavy optimization and those moving toward training-free or joint-modeling paradigms. While many frameworks, such as the Diffusion Transformer (DiT) that jointly models low-level VAE latents and high-level DINO semantic features, require unified training (\texttt{arXiv:2504.16064}), others prioritize inference-time flexibility. Group Diffusion introduces a training-free mechanism that allows Diffusion Transformers to share attention across image groups, improving ImageNet-256x256 FID by up to 32.2\% (\texttt{arXiv:2512.10954}). In the realm of single-image generation, a novel patch-based model replaces learned score functions with an analytical, closed-form denoiser derived from internal patch distributions (\texttt{arXiv:2606.04299}). This training-free approach offers extreme scalability—processing 1 GP in 13.9 minutes—and superior diversity (LPIPS Div. 0.50) compared to SinDiffusion (\texttt{arXiv:2606.04299}). However, its reliance on coarse-to-fine blending heuristics remains a weakness, as it must compensate for the lack of inherent global semantic modeling present in patch-based distributions (\texttt{arXiv:2606.04299}). Finally, to address the risk of memorization, Ambient Diffusion utilizes large noise scales during training to decouple generation from data memorization (\texttt{arXiv:2502.21278}).

\subsection{Tensions and Open Problems}

A fundamental tension persists regarding the architectural paradigm best suited for high-fidelity synthesis: the competition between adversarial frameworks and score-based diffusion models. Early advancements in stability and quality were driven by progressive growth strategies in GANs (\texttt{arXiv:1710.10196}) and style-based generator architectures (\texttt{arXiv:1812.04948}). However, recent critiques suggest that these adversarial methods struggle with the precise control and distributional coverage offered by diffusion. For instance, high-fidelity image editing frameworks (\texttt{arXiv:2312.15707}) argue that the inherent instabilities of GANs (\texttt{arXiv:1406.2661}) are superseded by the superior editing capabilities and structural reliability of diffusion-based approaches, such as GLIDE (\texttt{arXiv:2112.10741}).

A second debate centers on the mechanism of conditional guidance and its impact on sample fidelity. While classifier-free guidance (\texttt{arXiv:2207.12598}) has become a standard for aligning text-to-image models with linguistic prompts (\texttt{arXiv:2205.11487}), its limitations in fine-grained control have sparked significant pushback. Recent research (\texttt{arXiv:2312.13253}) critiques the reliance on classifier-free guidance, suggesting that pretrained generative models can achieve more precise conditional image generation through alternative conditioning mechanisms. This tension highlights a broader struggle in the field: whether to rely on implicit guidance through training-time alignment or to develop explicit, decoupled control mechanisms that do not compromise the underlying generative distribution.

Finally, the community is divided over the trade-off between sampling efficiency and generative complexity. The development of non-Markovian sampling paths, such as Denoising Diffusion Implicit Models (\texttt{arXiv:2010.02502}), aimed to accelerate generation, yet these methods face criticism regarding their ability to maintain long-term consistency during continual learning tasks (\texttt{arXiv:2311.14028}) or high-fidelity editing (\texttt{arXiv:2312.15707}). While some researchers advocate for complex, multi-step diffusion processes to ensure density coverage, others seek to bypass these iterative bottlenecks entirely through distillation or one-step consistency models (\texttt{arXiv:2303.01469}), creating a rift between those prioritizing mathematical rigor in the score-matching process and those prioritizing real-time deployment.

\subsection{Outlook}

The frontier of generative modeling is rapidly shifting from mere sample synthesis toward the explicit modeling of data manifold geometry (\texttt{arXiv:2606.00094}) and the pursuit of training-free, single-image generation (\texttt{arXiv:2606.04299}). As the field matures, the focus is moving away from brute-force scaling and toward structural efficiency, characterized by spectral progressive diffusion (\texttt{arXiv:2506.09376}) and interpretable diffusion trajectories via Koopman operators (\texttt{arXiv:2510.12220}). The next generation of models will likely bridge the gap between generative pre-training and one-step execution (\texttt{arXiv:2506.09376}), moving toward a unified framework where high-speed generation is no longer a compromise on quality, but a fundamental property of the learned manifold.

\section*{References}
\begin{enumerate}\small
  \item Auto-Encoding Variational Bayes (2013). \texttt{arXiv:1312.6114}
  \item Adam: A Method for Stochastic Optimization (2014). \texttt{arXiv:1412.6980}
  \item Deep Learning Face Attributes in the Wild (2014). \texttt{arXiv:1411.7766}
  \item Generative Adversarial Networks (2014). \texttt{arXiv:1406.2661}
  \item Microsoft COCO: Common Objects in Context (2014). \texttt{arXiv:1405.0312}
  \item LSUN: Construction of a Large-scale Image Dataset using Deep Learning
  with Humans in the Loop (2015). \texttt{arXiv:1506.03365}
  \item Deep Unsupervised Learning using Nonequilibrium Thermodynamics (2015). \texttt{arXiv:1503.03585}
  \item U-Net: Convolutional Networks for Biomedical Image Segmentation (2015). \texttt{arXiv:1505.04597}
  \item Improved Techniques for Training GANs (2016). \texttt{arXiv:1606.03498}
  \item GANs Trained by a Two Time-Scale Update Rule Converge to a Local Nash
  Equilibrium (2017). \texttt{arXiv:1706.08500}
  \item Progressive Growing of GANs for Improved Quality, Stability, and
  Variation (2017). \texttt{arXiv:1710.10196}
  \item Attention Is All You Need (2017). \texttt{arXiv:1706.03762}
  \item Decoupled Weight Decay Regularization (2017). \texttt{arXiv:1711.05101}
  \item Large Scale GAN Training for High Fidelity Natural Image Synthesis (2018). \texttt{arXiv:1809.11096}
  \item A Style-Based Generator Architecture for Generative Adversarial Networks (2018). \texttt{arXiv:1812.04948}
  \item Glow: Generative Flow with Invertible 1x1 Convolutions (2018). \texttt{arXiv:1807.03039}
  \item Neural Ordinary Differential Equations (2018). \texttt{arXiv:1806.07366}
  \item The Unreasonable Effectiveness of Deep Features as a Perceptual Metric (2018). \texttt{arXiv:1801.03924}
  \item Generative Modeling by Estimating Gradients of the Data Distribution (2019). \texttt{arXiv:1907.05600}
  \item Analyzing and Improving the Image Quality of StyleGAN (2019). \texttt{arXiv:1912.04958}
  \item Sliced Score Matching: A Scalable Approach to Density and Score
  Estimation (2019). \texttt{arXiv:1905.07088}
  \item Improved Precision and Recall Metric for Assessing Generative Models (2019). \texttt{arXiv:1904.06991}
  \item Denoising Diffusion Probabilistic Models (2020). \texttt{arXiv:2006.11239}
  \item Score-Based Generative Modeling through Stochastic Differential
  Equations (2020). \texttt{arXiv:2011.13456}
  \item Denoising Diffusion Implicit Models (2020). \texttt{arXiv:2010.02502}
  \item Improved Techniques for Training Score-Based Generative Models (2020). \texttt{arXiv:2006.09011}
  \item WaveGrad: Estimating Gradients for Waveform Generation (2020). \texttt{arXiv:2009.00713}
  \item DiffWave: A Versatile Diffusion Model for Audio Synthesis (2020). \texttt{arXiv:2009.09761}
  \item Training Generative Adversarial Networks with Limited Data (2020). \texttt{arXiv:2006.06676}
  \item Diffusion Models Beat GANs on Image Synthesis (2021). \texttt{arXiv:2105.05233}
  \item Improved Denoising Diffusion Probabilistic Models (2021). \texttt{arXiv:2102.09672}
  \item High-Resolution Image Synthesis with Latent Diffusion Models (2021). \texttt{arXiv:2112.10752}
  \item Variational Diffusion Models (2021). \texttt{arXiv:2107.00630}
  \item Image Super-Resolution via Iterative Refinement (2021). \texttt{arXiv:2104.07636}
  \item Cascaded Diffusion Models for High Fidelity Image Generation (2021). \texttt{arXiv:2106.15282}
  \item Knowledge Distillation in Iterative Generative Models for Improved
  Sampling Speed (2021). \texttt{arXiv:2101.02388}
  \item GLIDE: Towards Photorealistic Image Generation and Editing with
  Text-Guided Diffusion Models (2021). \texttt{arXiv:2112.10741}
  \item Learning Transferable Visual Models From Natural Language Supervision (2021). \texttt{arXiv:2103.00020}
  \item Score-based Generative Modeling in Latent Space (2021). \texttt{arXiv:2106.05931}
  \item Maximum Likelihood Training of Score-Based Diffusion Models (2021). \texttt{arXiv:2101.09258}
  \item Elucidating the Design Space of Diffusion-Based Generative Models (2022). \texttt{arXiv:2206.00364}
  \item Progressive Distillation for Fast Sampling of Diffusion Models (2022). \texttt{arXiv:2202.00512}
  \item Classifier-Free Diffusion Guidance (2022). \texttt{arXiv:2207.12598}
  \item Hierarchical Text-Conditional Image Generation with CLIP Latents (2022). \texttt{arXiv:2204.06125}
  \item DPM-Solver: A Fast ODE Solver for Diffusion Probabilistic Model Sampling
  in Around 10 Steps (2022). \texttt{arXiv:2206.00927}
  \item Photorealistic Text-to-Image Diffusion Models with Deep Language
  Understanding (2022). \texttt{arXiv:2205.11487}
  \item Flow Straight and Fast: Learning to Generate and Transfer Data with
  Rectified Flow (2022). \texttt{arXiv:2209.03003}
  \item On Distillation of Guided Diffusion Models (2022). \texttt{arXiv:2210.03142}
  \item Diffusion Model for Generative Image Denoising (2023). \texttt{arXiv:2302.02398}
  \item High-Fidelity Diffusion-based Image Editing (2023). \texttt{arXiv:2312.15707}
  \item Nested Diffusion Processes for Anytime Image Generation (2023). \texttt{arXiv:2305.19066}
  \item Consistency Models (2023). \texttt{arXiv:2303.01469}
  \item Diffusion Models Generate Images Like Painters: an Analytical Theory of
  Outline First, Details Later (2023). \texttt{arXiv:2303.02490}
  \item Simultaneous Image-to-Zero and Zero-to-Noise: Diffusion Models with
  Analytical Image Attenuation (2023). \texttt{arXiv:2306.13720}
  \item Simple diffusion: End-to-end diffusion for high resolution images (2023). \texttt{arXiv:2301.11093}
  \item Conditional Image Generation with Pretrained Generative Model (2023). \texttt{arXiv:2312.13253}
  \item Continual Learning of Diffusion Models with Generative Distillation (2023). \texttt{arXiv:2311.14028}
  \item Accelerated Image-Aware Generative Diffusion Modeling (2024). \texttt{arXiv:2408.08306}
  \item DiffuseHigh: Training-free Progressive High-Resolution Image Synthesis
  through Structure Guidance (2024). \texttt{arXiv:2406.18459}
  \item Diffusion Posterior Proximal Sampling for Image Restoration (2024). \texttt{arXiv:2402.16907}
  \item Enhancing Diffusion Models for High-Quality Image Generation (2024). \texttt{arXiv:2412.14422}
  \item Varying Manifolds in Diffusion: From Time-varying Geometries to Visual
  Saliency (2024). \texttt{arXiv:2406.18588}
  \item Generative Modeling with Diffusion (2024). \texttt{arXiv:2412.10948}
  \item Edge-preserving noise for diffusion models (2024). \texttt{arXiv:2410.01540}
  \item Diffusion Models in Low-Level Vision: A Survey (2024). \texttt{arXiv:2406.11138}
  \item Optical Diffusion Models for Image Generation (2024). \texttt{arXiv:2407.10897}
  \item Generative AI in Vision: A Survey on Models, Metrics and Applications (2024). \texttt{arXiv:2402.16369}
  \item Boosting Generative Image Modeling via Joint Image-Feature Synthesis (2025). \texttt{arXiv:2504.16064}
  \item Group Diffusion: Enhancing Image Generation by Unlocking Cross-Sample Collaboration (2025). \texttt{arXiv:2512.10954}
  \item MSF: Efficient Diffusion Model Via Multi-Scale Latent Factorize (2025). \texttt{arXiv:2501.13349}
  \item Does Generation Require Memorization? Creative Diffusion Models using Ambient Diffusion (2025). \texttt{arXiv:2502.21278}
  \item Enhancing diffusion models with Gaussianization preprocessing (2025). \texttt{arXiv:2512.21020}
  \item Hierarchical Koopman Diffusion: Fast Generation with Interpretable Diffusion Trajectory (2025). \texttt{arXiv:2510.12220}
  \item Denoising Score Distillation: From Noisy Diffusion Pretraining to One-Step High-Quality Generation (2025). \texttt{arXiv:2503.07578}
  \item Efficient Diffusion Training through Parallelization with Truncated Karhunen-Lo\textbackslash{}`eve Expansion (2025). \texttt{arXiv:2503.17657}
  \item Revisiting Diffusion Models: From Generative Pre-training to One-Step Generation (2025). \texttt{arXiv:2506.09376}
  \item Score-based Idempotent Distillation of Diffusion Models (2025). \texttt{arXiv:2509.21470}
  \item Diffusion Image Generation with Explicit Modeling of Data Manifold Geometry (2026). \texttt{arXiv:2606.00094}
  \item Efficient and Training-Free Single-Image Diffusion Models (2026). \texttt{arXiv:2606.04299}
  \item Spectral Progressive Diffusion for Efficient Image and Video Generation (2026). \texttt{arXiv:2605.18736}
\end{enumerate}

\end{document}
